\documentclass[../numerical_methods.tex]{subfiles}
\begin{document}
\section{Aula 11 - 14 de Abril, 2026}
\subsection{Motivações}
\begin{itemize}
	\item Métodos Iterativos para Sistemas Lineares
\end{itemize}
\subsection{Métodos Iterativos para Sistemas Lineares}
O primeiro tópico para além do tutorial de Python que iremos entrar é o uso de métodos numéricos diretos para resolução de sistemas lineares do tipo
\[
	Ax = b,
\]
onde A é uma matriz n por n e x e b são vetores.

Estes métodos são divididos em \textbf{métodos iterativos}, que fornecem uma sequência convergente de aproximações para o vetor solução a partir de um chute inicial \(\mathbf{x}^{(0)}\), e \textbf{métodos diretos}, que fornecem uma solução ``exata'' a menos de um erro de arredondamento, em um número finito de operações.

Para estudar este assunto, vamos revisar e introduzir conceitos básicos de sistemas lineares; o primeiro deles é a representação da equação dada, \(Ax=b\), em termos das colunas de A:
\[
	Ax = b \Longleftrightarrow a_1x_1+a_2x_2+\dotsc +a_{n}x_{n}=b,
\]
onde \(a_{j}\) é a j-ésima coluna de A. Desta forma, se A é \textit{não singular}, ou seja, se o determinante de A é não nulo, então as colunas \(a_{j}\) são linearmente independentes; logo, nestas condições, o vetor b pode ser unicamente escrito como combinação linear das colunas de A!

Em outras palavras, se A é uma matriz não singular, então o sistema \(Ax = b\) possui uma única solução. Algumas condições são equivalentes a isso, como:
\begin{itemize}
	\item[1)] \(\det{(A)}\neq 0\);
	\item[2)] As colunas ou linhas de A são linearmente independentes;
	\item[3)] Existe uma matriz inversa \(A^{-1}\) tal que \(AA^{-1}= A^{-1}A = I\);
	\item[4)] O conjunto gerado pelas colunas de A é um conjunto gerador de \(\mathbb{R}^{n}\); e
	\item[5)] O núcleo de A contém apenas o vetor nulo.
\end{itemize}

\begin{def*}
	Seja A uma matriz \(A\in \mathbb{M}_{n}(\mathbb{R})\); o \textbf{espaço coluna de A} é o conjunto \(C(A)\) contendo todas as combinações lineares das colunas de A. A dimensão deste espaço, \(\mathrm{dim}(C(A))\), é conhecido como \textbf{posto de A}, e o conjunto \(C(A^{T})\) é o \textbf{espaço linha de A}. \(\square\)
\end{def*}
É possível mostrar que \(\mathrm{dim}(C(A^{T})) = \mathrm{dim}(C(A))\).
\begin{def*}
	Seja A uma matriz \(A\in \mathbb{M}_{n}(\mathbb{R})\); o \textbf{espaço nulo de A}, denotado por \(N(A)\), é o que contém todos os vetores z tais que \(Az = 0\). \(\square\)
\end{def*}
Um resultado que relaciona os conceitos apresentados afirma que, dado que \(A\in \mathbb{M}_{n}(\mathbb{R})\), então
\[
	\mathrm{dim}(N(A))=n-\mathrm{dim}(C(A))
\]
\begin{example}
	Se considerarmos a matriz A como
	\[
		A = \begin{bmatrix}
			3 & 3 \\
			4 & 4 \\
		\end{bmatrix},
	\]
	então segue que
	\begin{align*}
		 & N(A)=\{\alpha (1,-1)^{T}:\; \alpha \in \mathbb{R}\}; \\
		 & C(A)=\{\beta (3,4)^{T}:\; \beta \in \mathbb{R}\}.
	\end{align*}
\end{example}

\begin{tcolorbox}[
		skin=enhanced,
		title=Observação,
		fonttitle=\bfseries,
		colframe=black,
		colbacktitle=cyan!75!white,
		colback=cyan!15,
		colbacklower=black,
		coltitle=black,
		drop fuzzy shadow,
		%drop large lifted shadow
	]

	Dado um sistema linear de ordem n, uma primeira ideia bem natural pra qualquer estudante que tenha passado por álgebra linear é utilizar a regra de Cramer para resolvê-lo:
	\[
		x_{i} = \frac{\det{(A_{i})}}{\det{(A)}},
	\]
	ou seja, não precisamos estudar métodos numéricos porque esse já funciona perfeitamente bem, não? Fim da disciplina!

	Mentira, se fosse assim, o mundo seria um lugar lindo e fácil de implementar sistemas lineares de qualquer ordem. No entanto, utilizando uma expansão de Laplace no cálculo do determinante, nota-se que são precisas \(3(n+1)!\) operações de pontos flutuantes para resolver um sistema assim à força bruta.
	Para dar uma noção disso, se \(n=20\), seriam necessários \(1,53 \times 10^{20}\) flops (Floating Point Operations) para resolver o sistema linear de tamanho 20, que nem é tão grande assim, mas um processador i7 de algumas gerações atrás realizava apenas \(9\times 10^{9}\) flops por segundo, resultando em apenas \(1,7\times 10^{10}\) segundos pra resolver, ou 539 anos.

	Dá pra ver que é inviável usar força bruta.
\end{tcolorbox}

Com matrizes triangulares superiores e inferiores, é mais plausível resolver sistemas lineares.
\begin{def*}
	Uma matriz \(L\in \mathbb{M}_{n}\) é dita \textbf{triangular inferior} se, para todo \(j > i \) (ou seja, se para todo par linha-coluna abaixo da diagonal principal), as entradas são nulas: \(\ell_{ij} = 0\). Por outro lado, uma matriz \(U\in \mathbb{M}_n\) é dita \textbf{triangular superior} se, para todo \(i>j\) (ou seja, se pra todo par linha-coluna acima da diagonal principal), as entradas são nulas: \(u_{ij}=0\). \(\square\)
\end{def*}

Assim, se A e B são duas matrizes triangulares inferior e superior respectivamente, qual seria a forma eficiente de calcular \(C = A \cdot B\)? Bom, neste caso,
\[
	c_{ij}=\sum\limits_{k=1}^{m}a_{ik}b_{kj}, \quad m = \min\limits_{}\{i, j\}.
\]
Daí, obtemos um sistema linear para cada uma delas; nas matrizes inferiores, ele será
\[
	\begin{bmatrix}
		a_{11} &        &        &        &        \\
		a_{21} & a_{22} &        &        &        \\
		a_{31} & a_{32} & a_{33} &        &        \\
		\vdots & \vdots & \vdots & \ddots &        \\
		a_{n1} & a_{n2} & a_{n2} & \dotsc & a_{nn}
	\end{bmatrix} \begin{bmatrix}
		x_{1}  \\
		x_{2}  \\
		x_{3}  \\
		\vdots \\
		x_{n}  \\
	\end{bmatrix} = \begin{bmatrix}
		b_{1}  \\
		b_{2}  \\
		b_{3}  \\
		\vdots \\
		b_{n}  \\
	\end{bmatrix},
\]
com \(a_{i}\neq 0\) para \(i=1,\dotsc , n.\) Daí, a solução de cada linha pode ser obtida por substituição direta:
\begin{align*}
	 & \tikz[baseline=-0.5ex]\draw[black, fill=black, radius=1.5pt](0, 0)circle;\; a_{11}x_1= b_1 \Rightarrow x_1 = \frac{b_1}{a_{11}}                                         \\
	 & \tikz[baseline=-0.5ex]\draw[black, fill=black, radius=1.5pt](0, 0)circle;\; a_{21}x_1 + a_{22}x_2=b_2 \Rightarrow x_{2}=\frac{b_{2}-a_{21}x_1}{a_{22}}                  \\
	 & \vdots                                                                                                                                                                  \\
	 & \boxed{\tikz[baseline=-0.5ex]\draw[black, fill=black, radius=1.5pt](0, 0)circle;\; x_{i}=\frac{b_{i}-\sum\limits_{j=1}^{i-1}a_{ij}x_{j}}{a_{ii}},\quad i=1, \dotsc , n}
\end{align*}

Nas matrizes superiores, caso tenha forma
\[
	\begin{bmatrix}
		a_{11} & a_{12} & a_{13} & \dotsc & a_{1n} \\
		       & a_{22} & a_{23} & \dotsc & a_{2n} \\
		       &        & a_{33} & \dotsc & a_{3n} \\
		       &        &        & \ddots & \vdots \\
		       &        &        &        & a_{nn}
	\end{bmatrix} \begin{bmatrix}
		x_{1}  \\
		x_{2}  \\
		x_{3}  \\
		\vdots \\
		x_{n}  \\
	\end{bmatrix} = \begin{bmatrix}
		b_{1}  \\
		b_{2}  \\
		b_{3}  \\
		\vdots \\
		b_{n}  \\
	\end{bmatrix},
\]
a solução também pode ser obtida por substituições diretas, de forma bem parecida:
\begin{align*}
	 & \tikz[baseline=-0.5ex]\draw[black, fill=black, radius=1.5pt](0, 0)circle;\; a_{nn}x_n= b_n \Rightarrow x_1 = \frac{b_n}{a_{nn}}                                                          \\
	 & \tikz[baseline=-0.5ex]\draw[black, fill=black, radius=1.5pt](0, 0)circle;\; a_{n-1, n-1}x_{n-1} + a_{n-1, n}x_{_n}=b_{_n-1} \Rightarrow x_{2}=\frac{b_{n-1}-a_{n-1, n}x_n}{a_{n-1, n-1}} \\
	 & \vdots                                                                                                                                                                                   \\
	 & \boxed{\tikz[baseline=-0.5ex]\draw[black, fill=black, radius=1.5pt](0, 0)circle;\; x_{i}=\frac{b_{i}-\sum\limits_{j=i+1}^{n}a_{ij}x_{j}}{a_{ii}},\quad i=1, \dotsc , n}
\end{align*}

Apesar de parecer mais complexo, este método já fornece uma melhoria na quantidade de flops necessários para resolver um sistema, precisando, para cada operação, de:
\begin{table}[H]
	\centering
	\caption{FLOPS para operações aritméticas básicas}
	\begin{tabular}{c | c | c | c}
		Soma & Subtração & Multiplicação & divisão \\
		\hline
		\hline
		i-2  & 1         & i-1           & 1       \\
	\end{tabular}
\end{table}
Assim,
\[
	\mathrm{Flops}=\sum\limits_{i=1}^{n}(i-2)+2\sum\limits_{i=1}^{n}1 + \sum\limits_{i=1}^{n}(i-1) = \frac{n(n-3)}{2} + 2n + \frac{n(n-1)}{2} = n^{2}.
\]

\begin{example}[Aplicando os Algoritmos de Subsittuição]
	Sabendo que uma matriz A tem decomposição LU como abaixo:
	\[
		\begin{bmatrix}
			2 & 0 & 1  \\
			4 & 3 & 7  \\
			6 & 6 & 16
		\end{bmatrix} = \begin{bmatrix}
			1 & 0 & 0 \\
			2 & 1 & 0 \\
			3 & 2 & 1
		\end{bmatrix} \cdot \begin{bmatrix}
			2 & 0 & 1 \\
			0 & 3 & 5 \\
			0 & 0 & 3
		\end{bmatrix},
	\]
	como resolvemos \(Ax = b\) no caso em que \(b = [2, 1, 4]^{T}\)?

	Segue que
	\[
		Ax = b \Longleftrightarrow (LU)x) = b \Longleftrightarrow L \cdot (Ux) = b.
	\]
	Chamando \(Ux\) de y, temos que resolver duas coisas: primeiro, resolver
	\[
		Ly = b
	\]
	com substituições progressivas; e, em seguida, resolver
	\[
		Ux = y
	\]
	com substituições regressivas. Tente implementar um código disso!
\end{example}

Tendo em vista o potencial da decomposição LU, dada uma matriz \(A\in \mathbb{M}_{n}\), queremos encontrar L e U em \(\mathbb{M}_{n}\) tal que
\[
	A = LU,
\]
onde
\[
	L = \begin{bmatrix}
		1         &           &           &        &   \\
		\ell_{21} & 1         &           &        &   \\
		\ell_{31} & \ell_{32} & 1         &        &   \\
		\vdots    & \vdots    & \vdots    & \ddots &   \\
		\ell_{n1} & \ell_{n2} & \ell_{n3} & \dotsc & 1
	\end{bmatrix} \quad\&\quad U = \begin{bmatrix}
		u_{11} & u_{12} & u_{13} & \dotsc & u_{1n} \\
		       & u_{22} & u_{23} & \dotsc & u_{2n} \\
		       &        & u_{33} & \dotsc & u_{3n} \\
		       &        &        & \ddots & \vdots \\
		       &        &        &        & u_{nn}
	\end{bmatrix}.
\]
Desta forma, o sistema ganha a cara
\[
	\begin{bmatrix}
		a_{11} & a_{12} & a_{13} & \dotsc & a_{1n} \\
		a_{21} & a_{22} & a_{23} & \dotsc & a_{2n} \\
		a_{31} & a_{32} & a_{33} & \dotsc & a_{3n} \\
		\vdots & \vdots & \vdots & \vdots & \vdots \\
		a_{n1} & a_{n2} & a_{n3} & \dotsc & a_{nn} \\
	\end{bmatrix} =\begin{bmatrix}
		1         &           &           &        &   \\
		\ell_{21} & 1         &           &        &   \\
		\ell_{31} & \ell_{32} & 1         &        &   \\
		\vdots    & \vdots    & \vdots    & \ddots &   \\
		\ell_{n1} & \ell_{n2} & \ell_{n3} & \dotsc & 1
	\end{bmatrix} \begin{bmatrix}
		u_{11} & u_{12} & u_{13} & \dotsc & u_{1n} \\
		       & u_{22} & u_{23} & \dotsc & u_{2n} \\
		       &        & u_{33} & \dotsc & u_{3n} \\
		       &        &        & \ddots & \vdots \\
		       &        &        &        & u_{nn}
	\end{bmatrix}.
\]

% TODO: CÓDIGO

Para a primeira linha de U, teremos \(u_{1j}=a_{1j}, \; j=1, \dotsc , n\); para a primeira coluna de L, será
\[
	\ell_{1j} = \frac{a_{i1}}{u_{11}},\; i = 1, 2, \dotsc , n;
\]
na segunda linha de U, teremos
\[
	u_{2j} = a_{2j} - \ell_{21}u_{1j}, \quad j=2, \dotsc , n;
\]
na segunda coluna de L,
\[
	\ell_{i2} = \frac{a_{i2} - \ell_{i1}u_{12}}{u_{22}},\quad i=3, \dotsc , n;
\]
e, por fim, o termo geral é:
\begin{itemize}
	\item Para \(u_{ij}\) com \(j\geq i\),
	      \begin{align*}
		      a_{ij} = \sum\limits_{k=1}^{i}\ell_{ik}u_{kj} & = \sum\limits_{k=1}^{i-1}\ell_{ik}u_{kj}+\ell_{ii}u_{ij}          \\
		      \Rightarrow                                   & \boxed{u_{ij} = a_{ij} - \sum\limits_{k=1}^{i-1}\ell_{ik}u_{kj}};
	      \end{align*}
	\item Para \(\ell_{ij}\) com \(i > j\),
	      \begin{align*}
		      a_{ij} = \sum\limits_{k=1}^{j} \ell_{ik}u_{kj} & =\sum\limits_{k=1}^{j-1}\ell_{ik}u_{kj} + \ell_{ij}u_{jj}                        \\
		      \Rightarrow                                    & \boxed{\ell_{ij} = \frac{a_{ij}-\sum\limits_{k=1}^{j-1}\ell_{ik}u_{kj}}{u_{jj}}}
	      \end{align*}
\end{itemize}

% TODO: CÓDIGO

\begin{def*}
	Seja \(A\in \mathbb{M}_{n}\). Os \textbf{menores principais de A} são as sub-matrizes da forma
	\[
		A_{k} = \begin{bmatrix}
			a_{11} & a_{12} & \dotsc & a_{1k} \\
			a_{21} & a_{22} & \dotsc & a_{2k} \\
			\vdots & \vdots & \ddots & \vdots \\
			a_{k1} & a_{k2} & \dotsc & a_{kk}
		\end{bmatrix}, \quad k=1, \dotsc , n.\; \square
	\]
\end{def*}
Com os menores principais, ganhamos um critério de existência e unicidade da decomposição LU!
\begin{theorem*}
	Sejam \(A\in \mathbb{M}_{n}\) e \(A_{k}\) seu menor principal de ordem k; se \(\det{(A_{k})}\neq 0\) para todo k entre 1 e n-1, então existe uma única L e uma única U, tais que
	\[
		A = LU.
	\]
\end{theorem*}
Mais ainda, ganhamos uma forma de calcular o determinante da matriz A: se \(A = LU\), então
\[
	\det{(A)} = \det{(LU)} = \det{(L)} \det{(U)};
\]
porém, \(\det{(L)} = 1\), ou seja,
\[
	\det{(A)} = \prod\limits_{i=1}^{n}u_{ii}
\]
\begin{exr}
	Seja
	\[
		A = \begin{bmatrix}
			1  & 2 & 0 \\ ]
			1  & 3 & 1 \\
			-2 & 0 & 1
		\end{bmatrix}
	\].
	% TODO: CÓDIGO

	Responda:
	\begin{itemize}
		\item[1)] A matriz A possui decomposição LU?
		\item[2)] Calcule \(A = LU\);
		\item[3)] Calcule \(\det{(A)}\) via \(LU\); e
		\item[4)] Resolva \(Ax = b\) com \(b = [3, 5, -1]^{T}\).
	\end{itemize}
\end{exr}


\end{document}
